\subsection{Finetuning}  %
\label{sec:t2v_post_training}

As in prior work~\citep{dai2023emu,emuvideo2023}, we improve the motion and aesthetic quality of the generated videos by finetuning the pre-trained model on a small finetuning set of selected videos.
The finetuning set videos and captions are manually curated, and thus we term this stage as supervised finetuning.
During this stage, we train multiple models and combine them to form the final model through a model averaging approach.
While our model can generate high quality images, we find that post-training specifically for images results in a significant boost in quality.
We describe the image-specific post-training recipe in~\Cref{sec:t2i_generation} and describe the video-specific post-training recipe next.











\noindent \textbf{Finetuning Video Data.}
We aim to collect a finetuning set of high quality videos with good motion, realness, aesthetics, with a wide range of concepts, and with high quality captions.
For finding such videos, we start with a large pool of videos and apply both automated and manual filtering steps (taking motivation from the curation recipe from~\citep{dai2023emu}).
There are four key stages that are run in sequence, each operating on the output of the previous stage:
(1) Establishing a set of candidate videos.
Here, we use automated filters that set strict thresholds on aesthetics, motion, scene change.
Additionally, we remove videos with small subjects using an object detection model~\citep{zhou2022detecting} on frames.
This stage results in a few million videos but with an unbalanced distribution of concepts.
(2) Balancing the concepts in set of videos.
The goal of this stage is to obtain a small enough subset of concept-balanced videos such that each can be manually filtered in the following steps.
We used our taxonomy of human verbs and expressions, defined in ~\cref{sec:pt-data}, to perform text $k$-NN methods
to retrieve videos for each concept from the candidate pool of videos.
We manually picked a few visually appealing seed videos per concept and performed video $k$-NN to get a concept-balanced subset of videos.
For $k$-NN, we used the video and text embeddings from a video-text joint embedding model.
(3) Manually identifying cinematic videos.
Many aspects to high quality finetuning data cannot be reliably captured by automated filters with high precision and recall.
At this stage, we instead rely on manual filtering.
Here we ensure that the remaining videos have angled (natural sunshine or studio) lighting, vivid (but not over-saturated) colors, no clutter, non-trivial motion, no camera shake, and no edited effects or overlay text.
During this stage, annotators additionally clip videos to the desired duration that will be trained on by selecting the best, most compelling clip of the video.
(4) Manually captioning the videos.
In detail, human annotators refine \llamaVideo generated captions by fixing incorrect details and ensuring the inclusion of certain key video details.
These include camera control, human expressions, subject and background info, detailed motion description and lighting information.
At this stage humans annotate six additional camera motion and position types (see~\Cref{appendix:data_camera_motion}).
Our video finetuning data is set to have duration between 10.6s and 16s.
In practice, 50\% of videos are 16s long, while the rest of 50\% videos are between 10.6s to 16s.

\noindent \textbf{Supervised Finetuning Recipe.} In video supervised finetuning (SFT), we use the same model architecture as the pre-training stage, and finetune the model with the pre-training checkpoints as initialization.
Different from pre-training that uses large-scale data, large batch sizes and training resources, we instead use relatively a small batch size and 64 nodes (512 H100 GPUs) to train the model, and use a cosine learning rate scheduler~\citep{loshchilov2016sgdr}.
Similar to the \pretraining stage, for videos that are at 16s, we train with 16 FPS, and for videos that are between 10.6s to 16s, we train with 24 FPS.
As a result, our model is trained to best support the generation of videos in both 10s and 16s.


\noindent \textbf{Model Averaging.}
Our experiments reveal that the choice of different sets of finetune data, hyperparameters as well as pre-train checkpoints significantly affects key aspects of the model's behavior, including motion, consistency, and camera control. To harness the diverse strengths of these models, we employ a model averaging approach.
Similar to \llama~\citep{llama3}, we average models obtained from SFT experiments that use various versions of finetune data, hyperparameters and pre-train checkpoints.





\subsection{Inference} \label{T2V_inference_details}
In this section, we describe the different hyper-parameters and settings used for sampling from \OursVideo.
For comparisons to prior work, we use a text classifier-free guidance scale of 7.5, and we use the linear-quadratic sampler described in~\cref{sec:linear_quad_sampler} with 50 steps (emulating 250 linear steps).
We also use an inference prompt rewrite on the input text prompt, as described below.
\subsubsection{Inference Prompt Rewrite} %
\label{sec:t2v_prompt_rewrite}

As mentioned in Section~\ref{sec:pt-data}, we train the model with high quality video/image-text pairs, and these training captions are characterized by their dense details and consistent paragraph structure.
However, the writing style and length of prompts in the inference stage vary widely.
For instance, most users typically type less than $10$ words, which is shorter than the average length of training captions.
To bridge the distribution gap between training captions and inference prompts, we utilize \llama~\citep{llama3} to transform the original input prompts into more detailed ones. The key details of the inference prompt rewrite model are:
\begin{itemize}
    \item We employ a standardized information architecture to rephrase the prompts, ensuring consistency in the visual composition.
    \item We refine the rewritten prompts by replacing complex vocabulary with more accessible and straightforward terminology, thereby enhancing their clarity and comprehensibility.
    \item We observe that excessively elaborate descriptions of motion details can result in the introduction of artifacts in the generated videos, highlighting the importance of striking a balance between descriptive richness and visual fidelity.
\end{itemize}



\textbf{Efficient Inference Rewrite Model.}
To improve the computation efficiency of the inference rewrite model, we developed a teacher-student distillation approach for this purpose.
Initially, we built a prompt rewrite teacher model based on the \llama 70B model, using detailed prompt instructions and in-context learning examples from the foundation model training set.
We then gathered human-in-the-loop (HITL) finetuning data.
This was achieved by using the \llama 70B prompt rewrite model as the teacher to conduct inference on a large prompt pool, and selecting high-quality rewrite pairs through human evaluations following the quality guideline.
Finally, we finetuned a 8B \llama model on the HITL prompt rewrite pairs to obtain the final prompt rewrite model to reduce the latency burden to the whole system.

